\section{Backward demands and finite cyclic closure}
\label{sec:demand}

Demand propagation is the bridge between a consumer-facing theorem and the producer precisions or resources
needed to prove it.  It should be a finite calculation, not an appeal to an unspecified continuity
principle.  This section gives exact algorithms for acyclic graphs and for monotone cyclic requirements
inside an explicit finite cap.

\subsection{Acyclic reverse propagation}

Let a finite computation graph have nodes \(0,\ldots,m-1\) in topological order.  Each edge
\(u\to v\) carries a primitive-recursive transformer
\[
  D_{u\leftarrow v}(q_v,x)
\]
which states how much of node \(u\)'s output is required to produce node \(v\)'s output at demand \(q_v\).
If node \(u\) has several consumers, combine their requests by a primitive-recursive join, usually maximum
coordinatewise.  Starting from the public demand at the final nodes, traverse the graph in reverse order:
\[
 q_u = q_u^{\mathrm{external}}\vee
       \bigvee_{u\to v}D_{u\leftarrow v}(q_v,x).
\]
Because the graph is finite and fixed, this reverse traversal is either unrolled into a concrete term or
implemented as a primitive recursion over a coded node table.  The output is a complete vector of demands.

\begin{theorem}[Acyclic demand soundness]
Suppose every local contract proves that inputs satisfying the edge demands yield the requested output
demand.  Then the reverse pass followed by the forward computation proves the final public contract.
\end{theorem}
\begin{proof}
Induct forward over the topological order.  At node \(v\), the reverse pass guarantees that every predecessor
was computed at least to \(D_{u\leftarrow v}(q_v,x)\).  Specialize the local continuity or correctness
contract at these demands, align its hypotheses with the already compiled predecessor contracts, and cut.
The staged graph theorem performs exactly this composition.  The final node has the public demand by
construction.
\end{proof}

For many applications this simple reverse pass is enough.  Addition of represented reals asks both inputs
for roughly the output precision.  Multiplication asks additionally for magnitude bounds.  Evaluation of a
truncated series propagates a coefficient cutoff and input precision to its coefficient and argument
producers.  A finite cover propagates the list of chart intersections that the final check will inspect.

\subsection{Why cycles arise}

Cycles appear even when the forward computation is acyclic.  An operation may need an input magnitude bound
to determine a precision, while that magnitude bound is itself certified at a precision depending on the
operation.  A branch-selection procedure may need a gap estimate, whose verifier uses a refined branch
approximation.  A local adaptive method may revise the requested radius after observing a finite error
certificate.

One must not simply say ``take a fixed point'' unless a finite constructive method is given.  The framework
uses a finite lattice with an explicit cap.

\subsection{Monotone closure in a finite box}

Fix \(r\) demand coordinates and caps \(M_0,\ldots,M_{r-1}\).  Let
\[
 L=\prod_{i<r}[0,M_i]
\]
with coordinatewise order.  A transformer \(F:L\to L\) is
\emph{inflationary} if \(q\le F(q)\), and monotone if \(q\le q'\) implies
\(F(q)\le F(q')\).  Starting from an initial request \(q^0\), define
\[
 q^{t+1}=F(q^t).
\]

\begin{theorem}[Bounded demand closure]
\label{thm:demand-closure}
If \(F\) is inflationary and maps the explicit finite box \(L\) into itself, then the iteration reaches a
fixed point after at most
\[
  1+\sum_{i<r}(M_i-q_i^0)
\]
steps.  The result is the least fixed point of \(F\) above \(q^0\).
\end{theorem}

\begin{proof}
As long as \(q^{t+1}\ne q^t\), inflationarity implies that at least one coordinate strictly increases, so
\(
 S(q)=\sum_i q_i
\)
increases by at least one.  It is bounded by \(\sum_i M_i\).  Hence a fixed point appears within the stated
fuel.  For leastness, let \(p\ge q^0\) be any fixed point.  By induction and monotonicity,
\(q^t\le p\) for every \(t\): the base is the assumption, and
\(q^{t+1}=F(q^t)\le F(p)=p\).  At stabilization, the computed fixed point is therefore below every such
\(p\).
\end{proof}

The execution is primitive recursive: the vector has fixed length, \(F\) is fixed, and the fuel is a
primitive-recursive sum.  A quantifier-free invariant records that every coordinate lies below its cap,
that the sequence is inflationary, and that the current state equals the \(t\)-fold iterate.  Leastness can
be exposed as a separate theorem parameterized by a candidate fixed point \(p\).

\subsection{Obtaining the cap}

The cap is not decorative.  It is a resource theorem.  There are several legitimate sources.

\begin{enumerate}
  \item A direct analytic bound: for example, every internal precision needed for output precision \(q\)
        is at most \(aq+b\) for fixed integers \(a,b\).
  \item A finite state space: demands are subsets of a fixed finite observation set, so the cap is its
        cardinality or bit mask.
  \item A supplied boundary datum: a condition number, separation bound, coefficient height, or finite
        cover determines the cap.
  \item A syntactic graph bound: a fixed number of nodes and local transformers with known affine growth
        yield a primitive-recursive global upper bound by repeated composition.
\end{enumerate}

Without a cap, monotone iteration over \(\Nat^r\) need not stabilize.  The framework therefore classifies a
claim of cyclic demand closure without a cap as incomplete.

\subsection{Noninflationary transformers}

A raw system of equations may give a monotone but not inflationary transformer \(G\).  Replace it by
\[
 F(q)=q\vee G(q).
\]
Fixed points of \(F\) are post-fixed points of \(G\), which is exactly what demand soundness requires: the
allocated resource dominates every locally requested resource.  The least \(F\)-fixed point above the
external request is the least closed allocation.  This transformation is primitive recursive and makes the
termination argument transparent.

\subsection{Demand closure as a loop node}

The closure computation becomes a local loop contract:
\begin{itemize}
  \item precondition: the initial request and cap are valid, and \(F\) maps the box into itself;
  \item state: \((t,q^t,\mathit{stable})\);
  \item invariant: validity, box membership, inflationarity along the trace, and correctness of the iterate;
  \item fuel: \(1+\sum_i(M_i-q_i^0)\);
  \item postcondition: \(F(q)=q\) and \(q\) dominates the initial request.
\end{itemize}
The loop compiles by Theorem~\ref{thm:loop-compile}, and the resulting implication is inserted as a node in
the static proof graph.  The forward computation then uses the closed demands as ordinary parameters.

\subsection{Backward demand for changing-parameter recursion}

Consider a forward recurrence
\[
 s_{i+1}=T(i,s_i,p_i),
 \qquad p_{i+1}=U(i,s_i,p_i),
\]
and suppose a request on the final state at time \(N\) induces requests on earlier states and parameters.
For a fixed horizon, compute demands backward:
\[
 d_N=q,
 \qquad d_i=B(i,d_{i+1},\widehat s_i,\widehat p_i),
\]
where \(\widehat s_i,\widehat p_i\) are coarse envelopes known before the precise run.  Then replay forward
at the allocated demands.  The proof has two layers:
\begin{enumerate}
  \item a reverse finite induction or trace proves that each \(d_i\) is sufficient for the next local step;
  \item a forward quantifier-free induction proves that the actual approximations meet those demands.
\end{enumerate}
If the backward transformer itself needs information refined by the forward run, enclose both in the finite
cyclic closure above.  The cap must include both precision and envelope coordinates.

\subsection{Metastable consumers}

Some theorems do not need a global convergence modulus.  A finite consumer inspects only a window
\([N,N+g(N)]\) and asks for approximate constancy there.  The demand is then the finite function code for
\(g\), an error \(q\), and a search bound for \(N\).  A bounded search through candidate windows, with a
Boolean verifier for all pairs in the window, is a deterministic finite contract once the search bound is
available.  This form is often substantially weaker than a Cauchy modulus and is compatible with iPRA.
The framework keeps the two statements distinct: a metastable witness does not silently become a full
modulus.

\subsection{A failure mode exposed by the framework}

Suppose a draft defines demands by equations
\[
 q_1=q+q_2,
 \qquad q_2=q_1+1.
\]
No finite solution exists.  A vague instruction to ``iterate until stable'' hides the contradiction.  The
finite-box method catches it: either the transformer escapes every proposed cap, or it reaches the cap
without a fixed point.  Likewise, a discontinuous branch selector may create an unbounded demand near a
collision.  The correct repair is not a larger arbitrary cap but a relational output or an explicit gap
assumption.  Demand analysis is therefore also a diagnostic tool.
